% ============================================================================
% Chapter 2: Literature Review & Background
% ============================================================================
\chapter{Background and Related Work}
\label{ch:literature}

This chapter reviews the research landscape relevant to CBSR across three areas: provenance-based
intrusion detection, causal anomaly detection methods, and conformal prediction. We conclude with
a research gap analysis that motivates our contributions.

\section{Provenance-Based Intrusion Detection Systems}

Provenance graphs represent system history by capturing relationships between system entities
(processes, files, network sockets) as directed edges. Each edge encodes an operation such as
\texttt{PROCESS:chrome.exe $\to$ READS\_FILE $\to$ FILE:History}, preserving the complete causal
lineage of system execution.

\subsection{Supervised and Semi-Supervised Approaches}

\textbf{OCR-APT}~\cite{ocrapt2025} reconstructs APT attack stories from audit logs using subgraph
anomaly detection combined with LLMs. It reports AUC of 0.91 on DARPA Trace logs but requires
24--48 hours of clean pre-collection to establish a benign baseline.

\textbf{TraceCluster}~\cite{tracecluster2026} proposes a lightweight clustering-based subgraph
attention network for APT detection, achieving average precision of 0.88 on multi-hop APT
scenarios. Its GNN is trained on labeled subgraph clusters---requiring labeled examples.

\textbf{StageFinder}~\cite{stagefinder2026} learns temporal reasoning over provenance data for
attack stage estimation, reporting 94\% recall on staged APT campaigns but requiring labeled
attack stage annotations.

\subsection{Unsupervised and Reduced-Label Approaches}

\textbf{UNICORN}~\cite{han2020unicorn} (NDSS 2020) summarizes streaming provenance subgraphs as
graph histograms and fits an unsupervised clustering baseline. It reduces the label dependency but
still requires a \emph{clean, attack-free} training period to define normal cluster centroids.

\textbf{WATSON}~\cite{gao2020watson} (CCS 2020) aggregates audit events using contextual
semantics but relies on static, clean offline profiles with no mechanism to adapt to concept drift.

\textbf{ShadeWatcher}~\cite{yang2022shadewatcher} (S\&P 2022) frames threat analysis as a
recommendation task. It requires a clean audit log baseline and leverages user feedback (labels)
for guidance.

\textbf{Kairos}~\cite{cheng2024kairos} utilizes self-supervised learning over provenance graphs
to minimize label dependency, yet still depends on an uncontaminated pre-collection phase for
initial parameter fitting.

\subsection{Critical Assessment}

Table~\ref{tab:lit_requirements} summarizes the deployment requirements of existing
provenance-based IDS. A clear pattern emerges: \emph{every existing system} requires some form of
clean baseline data. This shared assumption creates a systemic vulnerability.

\begin{table}[htbp]
\centering
\small
\caption{Capability comparison of provenance IDS methods. ``Clean Baseline'' requires a clean,
attack-free pre-collection period; ``FPR Guarantee'' provides a formal bound on false alarm rate.}
\begin{tabularx}{\linewidth}{Xcccc}
\hline
\rowcolor[gray]{0.88}
\textbf{Method} & \textbf{Clean Base.} & \textbf{Drift Adapt.} & \textbf{Causal} & \textbf{FPR Guar.} \\
\hline
OCR-APT~\cite{ocrapt2025}       & Yes (48h) & No  & No  & No  \\
TraceCluster~\cite{tracecluster2026} & Yes  & No  & No  & No  \\
StageFinder~\cite{stagefinder2026}   & Yes  & No  & No  & No  \\
Kairos~\cite{cheng2024kairos}        & Yes  & No  & No  & No  \\
UNICORN~\cite{han2020unicorn}         & Yes  & No  & No  & No  \\
Causal-IDS~\cite{causalids2026}       & Yes  & No  & Yes & No  \\
\textbf{CBSR (Ours)}                  & \textbf{No} & \textbf{Yes} & \textbf{Yes} & \textbf{Yes} \\
\hline
\end{tabularx}
\label{tab:lit_requirements}
\end{table}


\section{Causal Anomaly Detection}

Causal reasoning is increasingly applied to cybersecurity, motivated by the insight that causal
relationships are more robust to distributional shifts than statistical correlations.

\textbf{Causal-IDS}~\cite{causalids2026} (ICOIN 2026) is the closest prior work. It models
network flow log variables using static Structural Causal Models (SCMs) to identify intrusions as
violations of learned causal mechanisms. However, it differs from CBSR in three fundamental ways:
(1) it operates on coarse-grained network flow statistics rather than fine-grained provenance
edges; (2) it requires clean training logs; and (3) it uses fixed empirical thresholds with no
drift adaptation or FPR control.

\textbf{CausalGraph}~\cite{causalgraph2026} (ICCE 2026) integrates causal reasoning with LLMs for
intrusion detection. While offering superior interpretability, the computational cost of LLM
inference makes it impractical for real-time, high-velocity streaming logs.

\textbf{LiNGAM-SF}~\cite{lingamsf2025} (IEEE TNNLS 2025) proposes causal structural learning with
Linear Non-Gaussian Acyclic Models for streaming features. Its key insight---distinguishing latent
confounders from direct mechanism violations---informs the design of CBSR's behavioral state
coordinates.


\section{Conformal Prediction}

Conformal prediction~\cite{shafer2008, vovk2005algorithmic} provides a distribution-free
framework for constructing prediction sets with guaranteed coverage probabilities. Under the sole
assumption of exchangeability, the conformal p-value for a new test score $S_t$ against a
calibration set $\{s_1, \ldots, s_M\}$ is:
\begin{equation}
\text{conf\_pval}(S_t) = \frac{|\{i : s_i \geq S_t\}| + 1}{M + 1}
\end{equation}
This guarantees $\Pr(\text{false alarm}) \leq \alpha$ for any single test point. In streaming
security telemetry, exchangeability is violated by concept drift. CBSR adopts online threshold
adaptation~\cite{barber2023conformal, angelopoulos2021gentle} combined with change-point-triggered
queue flushing to restore localized exchangeability after drift events.


\section{Research Gap Analysis}

Based on the literature review, we identify four critical gaps that no existing system addresses
simultaneously:

\begin{enumerate}
\item \textbf{Semi-supervised causal IDS for provenance graphs}: No existing system combines causal
    mechanism violation detection with provenance-graph analysis using only sparse validation labels
    without requiring clean training data.
\item \textbf{Online SCM adaptation}: Existing causal IDS use static SCMs with no mechanism for
    concept drift recovery.
\item \textbf{Provable false alarm control}: No provenance-based IDS provides statistical
    guarantees on false alarm rates.
\item \textbf{Multi-perspective behavioral state reasoning}: No prior work fuses causal,
    temporal, and behavioral perspectives on a shared formal state representation.
\end{enumerate}

CBSR addresses all four gaps simultaneously.
